\section{Conclusions}

This paper evaluated whether national high-resolution hydrography and environmental datasets can be converted into structurally consistent, executable, and auditable SWAT\textsuperscript{+} projects through a documented automated workflow implemented in SWATGenX.

SWATGenX demonstrates that NHDPlus HR and national ancillary data can be converted automatically into executable, auditable SWAT\textsuperscript{+} project packages with documented preprocessing trade-offs and evidence that those packages can enter standard downstream calibration, sensitivity, and uncertainty workflows.
Primary evidence comprises a layered implementation with inspectable outputs and a stated preprocessing rule set that yields quantified structural changes relative to original NHDPlus HR sources, together with a drainage-area fidelity audit showing that executable channel areas preserve original NHDPlus HR cumulative drainage at gaged channels (85\% of matched gages within $0.5$--$2.0\times$).
Secondary evidence comprises structural and assembly-cost scaling across three benchmark size classes and workflow-readiness tests on two controlled gage basins.
Supporting evidence comprises measured simulation cost on one documented host/build configuration.

The evaluation does not establish nationwide hydrologic validity, comparative superiority to other national platforms, or readiness for regulatory filing without expert review.
Further work should isolate hydrologic consequences of alternative preprocessing choices and expand domain coverage beyond the selected evaluation matrix (Table~\ref{tab:model-roster}).
